\begin{tikzpicture}[
transform shape,
node distance=0.5cm and 1.25cm, % Vertical and horizontal distance
boxnode/.style={draw, rectangle, minimum height=0.8cm, text width=2.5cm, align=center, font=\small}, % Style for generic boxes
textnode/.style={text width=2.1cm, align=center, font=\small}, % Style for descriptive text
measurenode/.style={text width=2.5cm, align=center, font=\small}, % Style for measure names
mathnode/.style={font=\large}, % Larger font for main math symbols
tinytext/.style={font=\tiny, align=left} % For the example of non-measurable set
]

\node (E) [mathnode] {环 $\mathscr{E}$};
\node (RE) [mathnode, right=of E] {$\mathscr{R}_{\sigma}(\mathscr{E})$};
\node (M) [mathnode, right=of RE] {$\mathscr{M}$};
% \node (SE) [mathnode, right=of M] {$\overset{\scriptstyle\underset{\mkern4mu\verteq}{\{ E ~:~ E \subset \bigcup\limits_{n=1}^{\infty} A_n, A_n \in \mathscr{E} \}}}{\mathscr{S}(\mathscr{E})} = \mathscr{P}(\mathbb{R}^n)$};
\node (SE) [mathnode, right=of M] {$\mathscr{S}(\mathscr{E}) = \mathscr{P}(\mathbb{R}^n)$};
\node (verteq) [textnode, above=of SE, yshift=-1.7em, xshift=-0.8cm] {$\verteq$};
\node [above=of verteq, yshift=-2.1em, xshift=0.75cm] {\smaller[1] $\{ E ~:~ E \subset \bigcup\limits_{n=1}^{\infty} A_n, A_n \in \mathscr{E} \}$};

\node (subset1) at ($(E.east)!0.5!(RE.west)$) {$\subsetneqq$};
\node (subset2) at ($(RE.east)!0.5!(M.west)$) {$\subsetneqq$};
\node (subset3) at ($(M.east)!0.5!(SE.west)$) {$\subsetneqq$};

\node (basis) [measurenode, below=of E, yshift=0.5cm] {\smaller[1] 半开闭矩体有限并, 及 $\emptyset$};
\node (borel_set) [measurenode, below=of RE, yshift=0.5cm] {\smaller[1] Borel 集类};
\node (sigma_gen) [textnode, above=of RE, yshift=-1.2em] {\smaller[1] 由 $\mathcal{E}$ 生成的 $\sigma$ 环};
\node (non_borel) [textnode, above=of subset2, yshift=0.4cm] {存在非 Borel 集的可测集};
\node (lebesgue_measure) [measurenode, below=of M, yshift=0.5cm] {\smaller[1] Lebesgue 可测集类};
\node (non_meas) [textnode, above=of subset3, yshift=0.4cm] {存在不可测集};
\node (all_subsets) [measurenode, below=of SE, yshift=0.5cm] {\smaller[1] $\mathbb{R}^n$ 全体子集};

\node (mu1) [mathnode, below=of E, yshift=-2em] {$\mu$};
\node (tilde_mu1) [mathnode, below=of RE, yshift=-2em] {$\tilde{\mu}$};
\node (tilde_mu2) [mathnode, below=of M, yshift=-2em] {$\tilde{\mu}$};
\node (mu_star) [mathnode, below=of SE, yshift=-2em] {
$\mu^* E = \inf\limits_{E \subset \bigcup A_n} \sum \mu(A_n)$
};
\node (mu_final) [mathnode, below=of mu_star, yshift=-2cm] {};

\draw[->] (non_borel) -- (subset2);
\draw[->] (non_meas) -- (subset3);

\draw[->] (tilde_mu1) to [bend right] (mu1);
\draw[->] (tilde_mu2) to [bend right] (tilde_mu1);
\draw[->] (mu_star) to [bend right] (tilde_mu2);

% Long curved arrow
\draw[->] (mu1.south) to [bend right] (mu_star.south west); % Curved arrow from mu1 to mu_star

\end{tikzpicture}
